% Part: first-order-logic
% Chapter: model-theory
% Section: nonstandard-arithmetic

\documentclass[../../../include/open-logic-section]{subfiles}

\begin{document}

\olfileid{mod}{bas}{nsa}
\section{અંકગણિતના અપ્રમાણભૂત નિદર્શો}

\begin{defn}
$\Lang{L_N}$ એ અંકગણિતની !!{language} હોય,
જેમાં !!{constant}~$\Obj{0}$, દ્વિસ્થાની
!!{predicate}~$<$, એકસ્થાની !!{function}~$\prime$
અને દ્વિસ્થાની !!{function}s~$+$ તથા $\times$
છે.
\begin{enumerate}
\item અંકગણિતનું \emph{પ્રમાણભૂત નિદર્શ}
  એ $\Lang{L_N}$ માટેની !!{structure}~$\Struct{N}$
  છે, જેનું વ્યાપકક્ષેત્ર
  $\Nat = \{ 0, 1, 2, \dots\}$ છે અને જે
  $\Obj{0}$નું અર્થઘટન $0$ તરીકે, $<$નું
  $\Nat$ પરના નાનાપણાના સંબંધ તરીકે,
  તથા $\prime$, $+$ અને $\times$નું
  અનુક્રમે $\Nat$ પર ઉત્તરગામી,
  સરવાળો અને ગુણાકાર તરીકે કરે છે.
\item \emph{સાચું અંકગણિત} એ
  સિદ્ધાંત $\Theory{N}$ છે.
\end{enumerate}
\end{defn}

$\Lang{L_N}$માં કામ કરતાં, $\Obj{0}$ પર
ઉત્તરગામી વિધેય $n$ વાર લગાવી મળતા
$\Obj{0}^{\prime\dots\prime}$
સ્વરૂપના દરેક પદને ટૂંકમાં~$\num{n}$
લખીએ છીએ.

\begin{defn}
$\Lang{L_N}$ માટેની !!{structure}~$\Struct{M}$
\emph{પ્રમાણભૂત} ત્યારે અને માત્ર ત્યારે
જ જ્યારે $\Struct{N} \iso \Struct{M}$.
\end{defn}

\begin{thm}\ollabel{thm:non-std}
સાચા અંકગણિતના અપ્રમાણભૂત
!!{enumerable} નિદર્શો અસ્તિત્વમાં છે.
\end{thm}

\begin{proof}
$\Lang{L_N}$માં નવું !!{constant}~$c$
ઉમેરી તેને વિસ્તારો અને આ સિદ્ધાંત વિચારો:
\[
\Theory{N} \cup \Setabs{\num{n} < c}{n \in \Nat}.
\]
આ સિદ્ધાંત સાન્ત રીતે સંતોષ્ય છે.
તેથી સઘનતા મુજબ તેને નિદર્શ~$\Struct{M}$
છે. અધોગામી લેવેનહાઇમ--સ્કોલેમ પ્રમેય
મુજબ એ નિદર્શ !!{enumerable} પસંદ
કરી શકાય છે. $\Domain{M}$ એ
$\Struct{M}$નું વ્યાપકક્ષેત્ર હોય.
$\Struct{M}$માં $\Lang{L_N}$ના
અતાર્કિક સંકેતોનાં અર્થઘટનને આ રીતે
નામ આપીએ: $\mathbf{z} = \Assign{\Obj{0}}{M} \in \Domain{M}$,
${\prec} = \Assign{<}{M} \subseteq M^2$,
$* = \Assign{\prime}{M} \colon
\Domain{M} \to \Domain{M}$, અને
$\oplus = \Assign{+}{M}, \otimes =
\Assign{\times}{M} \colon \Domain{M}^2 \to \Domain{M}$.
દરેક $x \in \Domain{M}$ માટે,
$x$ પર~$*$ લગાવી મળતા
$\Domain{M}$ના ઘટકને $x^*$
લખીએ છીએ.
\sourcecorrection{OLMD-190}{મૂળ સાબિતીના
આ વાક્યમાં અચાનક $\Lang{L}$ લખાયું છે;
વિસ્તારેલી ભાષાના મૂળ અંકગણિતીય
સંકેતો $\Lang{L_N}$ના છે, તેથી
અહીં એ જ નામ પુનઃસ્થાપિત કર્યું છે.}

હવે જો $h$ એ $\Struct{N}$ અને
$\Struct{M}$ વચ્ચેની એકરૂપતા હોત,
તો કોઈ $n \in \Nat$ માટે
$h(n) = \Assign{c}{M}$ થાત.
એટલે $\Struct{N}$માં એવી
કોઈ પણ નિયુક્તિ~$s$ લો કે
$s(x) = n$. ત્યારે
$\Sat{N}{\eq[\num{n}][x]}[s]$.
\olref[iso]{thm:isom}ની સાબિતી
મુજબ $\Sat{M}{\eq[\num{n}][x]}[h \circ s]$
પણ છે; તેથી $\Assign{c}{M} =
\mathbf{z}^{*\cdots *}$, જ્યાં
$*$ને $n$ વાર લગાવ્યું છે.
પરંતુ ધારણા મુજબ
$\Sat{M}{\num{n} < c}$ અને
$\prec$ અસ્વવાચક હોવાથી
આ અશક્ય છે. તેથી
$\Struct{M}$ અપ્રમાણભૂત છે.
\end{proof}

\begin{prob}
ગણ~$X$ પરનો સંબંધ $R$
\emph{સુઆધારિત} ત્યારે અને માત્ર ત્યારે
જ જ્યારે $R$માં અનંત અધોગામી
સાંકળ ન હોય; એટલે $X$માં
એવાં $x_0$, $x_1$, $x_2$,~\dots
ન હોય કે $\dots x_2Rx_1Rx_0$.
ઝેર્મેલો--ફ્રેન્કેલ ગણસિદ્ધાંત~$ZF$
સુસંગત છે એમ માની, $ZF$ના
સુઆધારિત ન હોય એવા નિદર્શો
અસ્તિત્વમાં છે તે બતાવો:
એવાં નિદર્શો $\mathfrak{M}$
કે $\dots x_2 \in x_1 \in x_0$.
\end{prob}

\olref{thm:non-std}માંથી મળેલું
અપ્રમાણભૂત નિદર્શ $\Struct{M}$
પ્રમાણભૂત નિદર્શ સાથે
પ્રાથમિક રીતે સમકક્ષ હોવાથી
$\Struct{M}$ના અનેક ગુણધર્મો
મેળવી શકાય છે. આ વિભાગનો
બાકીનો ભાગ એ જ કામને સમર્પિત
છે; તેના વડે $\Theory{N}$ના
!!{enumerable} અપ્રમાણભૂત
નિદર્શોનું ચોક્કસ લક્ષણ
મળશે.

\begin{enumerate}
\item $\Domain{M}$નો કોઈ ઘટક
  પોતાનાથી $\prec$-નાનો નથી:
  $\lforall[x][\lnot x < x]$
  વાક્ય $\Struct{N}$માં સાચું
  છે અને તેથી $\Struct{M}$માં પણ.
\item એ જ પ્રકારની દલીલથી
  $\prec$ એ $\Domain{M}$નો
  \emph{રેખીય ક્રમ} છે; એટલે
  $\Domain{M}$ પરનો સંપૂર્ણ,
  અસ્વવાચક અને પરંપરિત સંબંધ.
\item $\mathbf{z}$ એ
  $\Domain{M}$નો $\prec$-લઘુતમ
  ઘટક છે.
\item $\Domain{M}$નો દરેક
  ઘટક તેના $*$-ઉત્તરગામી
  કરતાં $\prec$-નાનો છે,
  અને $x^*$ એ $x$ કરતાં
  મોટો $\Domain{M}$નો
  $\prec$-લઘુતમ ઘટક છે.
\item $\Struct{M}$માં
  $\Nat$ને એકરૂપ એવો
  $\prec$નો આરંભિક ખંડ છે:
  $\mathbf{z}, \mathbf{z}^*, \mathbf{z}^{**}, \dots$.
  તેને $\Domain{M}$નો
  \emph{પ્રમાણભૂત ભાગ}
  કહીએ છીએ. $\Domain{M}$નો
  અન્ય દરેક ઘટક
  \emph{અપ્રમાણભૂત} છે.
  $\Domain{M}$માં અપ્રમાણભૂત
  ઘટકો હોવા જ
  જોઈએ; નહિતર
  \olref{thm:non-std}ની
  સાબિતીનું વિધેય~$h$
  એકરૂપતા હોત.
  $\Struct{M}$ના આ પ્રમાણભૂત
  ભાગમાં ફરતા !!{variable}s
  માટે $n, m, \dots$
  વાપરીશું.
\item દરેક અપ્રમાણભૂત
  ઘટક દરેક પ્રમાણભૂત
  ઘટક કરતાં મોટો છે.
  કારણ કે દરેક
  $n \in \Nat$ માટે
  \[
  \Sat{N}{\lforall[z][(\lnot(\eq[z][\Obj{0}] \lor
  \dots \lor \eq[z][\num{n}]) \lif \num{n} < z)]},
  \]
  તેથી $z \in \Domain{M}$
  બધા પ્રમાણભૂત ઘટકોથી
  જુદો હોય તો તે
  એમના બધા કરતાં
  \emph{મોટો} હોવો જોઈએ.
\item $\mathbf{z}$ સિવાય
  $\Domain{M}$નો દરેક
  ઘટક કોઈ એકમાત્ર
  $\Domain{M}$-ઘટકનો
  $*$-ઉત્તરગામી છે;
  તેને $^*x$થી દર્શાવીએ.
  જો $x = y^*$ હોય,
  તો $x$ અને $y$
  બંનેમાંથી એક
  પ્રમાણભૂત હોય ત્યારે
  બંને પ્રમાણભૂત છે;
  અને એક અપ્રમાણભૂત
  હોય ત્યારે બંને
  અપ્રમાણભૂત છે.
\item $\Domain{M}$ પર
  સમકક્ષતા-સંબંધ
  $\approx$ આમ
  વ્યાખ્યાયિત કરો:
  કોઈ \emph{પ્રમાણભૂત}
  $n$ માટે $x \oplus n = y$
  અથવા $y \oplus n =x$
  હોય ત્યારે અને માત્ર
  ત્યારે $x \approx y$.
  બીજા શબ્દોમાં,
  $x$ અને $y$ વચ્ચે
  સાન્ત અંતર હોય
  ત્યારે અને માત્ર
  ત્યારે $x \approx y$.
  $n$ અને $m$
  પ્રમાણભૂત હોય તો
  $n \approx m$.
  $x$નો \emph{ખંડ}
  એટલે તેનો સમકક્ષતા-વર્ગ
  $[x] = \Setabs{y}{x
  \approx y}$.
\item ધારો કે $x \prec y$
  અને $x \not\approx y$.
  $\Sat{N}{\lforall[x][\lforall[y][(x < y \lif (x' < y \lor x' =
      y))]]}$ હોવાથી કાં તો
  $x^* \prec y$ અથવા
  $x^* = y$. બીજું
  અશક્ય છે, કારણ કે
  તેનાથી $x \approx y$
  થાય; તેથી $x^* \prec
  y$. એ જ રીતે,
  જો $x \prec y$
  અને $x \not\approx y$
  હોય, તો $x \prec
  {^*y}$. તેથી
  $x \prec y$ અને
  $x \not\approx y$
  હોય તો દરેક
  $w \approx x$ એ
  દરેક $v \approx y$
  કરતાં $\prec$-નાનો છે.
  તદનુસાર, દરેક
  અપ્રમાણભૂત ખંડ~$[x]$
  નીચે જેવી બંને
  દિશાએ અનંત સાંકળ
  બનાવે છે:
  \[
  \cdots \prec  {^{**}x} \prec {^*}x \prec x \prec x^* \prec x^{**}
  \prec \cdots
  \]
  પૂર્ણાંકો જેવો ક્રમપ્રકાર
  હોવાથી તેને $Z$-સાંકળ
  કહે છે.
  \sourcecorrection{OLMD-191}{મૂળમાં
  $x^*=y$ અશક્ય ઠરાવ્યા પછી
  ભૂલથી ફરી $x\prec y$ લખાયું છે;
  સાચો નવો નિષ્કર્ષ
  $x^*\prec y$ છે.}
  \sourcecorrection{OLMD-192}{મૂળ
  દરેક ખંડને બંને દિશાએ
  અનંત $Z$-સાંકળ કહે છે,
  પરંતુ પ્રમાણભૂત ખંડમાં
  લઘુતમ $\mathbf z$ છે.
  આ વર્ણન માત્ર
  અપ્રમાણભૂત ખંડો માટે
  સાચું છે.}
\item $\prec$-ક્રમ ખંડો
  પર પણ મૂકી શકાય છે:
  $x \prec y$ હોય અને
  બંને જુદા ખંડોમાં
  હોય, તો $x$નો
  ખંડ $y$ના ખંડ
  કરતાં નાનો છે.
  અપ્રમાણભૂત ઘટક
  ધરાવતો ખંડ
  \emph{અપ્રમાણભૂત}
  છે. પ્રમાણભૂત
  ખંડ લઘુતમ છે.
  \sourcecorrection{OLMD-193}{મૂળ
  $x\prec y$ માત્રથી
  $[x]$ એ $[y]$
  કરતાં નાનો કહે છે;
  એક જ ખંડનાં
  જુદાં ઘટકો માટે
  એ ખોટું છે. અહીં
  જુદા ખંડોની શરત
  સ્પષ્ટ કરી છે.}
\item લઘુતમ અપ્રમાણભૂત
  ખંડ નથી. જો $y$
  અપ્રમાણભૂત હોય,
  તો એવો $x \prec y$
  છે કે $x$ પણ
  અપ્રમાણભૂત છે અને
  $x \not\approx y$.
  સાબિતી: પ્રમાણભૂત
  નિદર્શ $\Struct{N}$માં
  દરેક સંખ્યાને
  બે વડે ભાગી શકાય,
  કદાચ એક શેષ સાથે:
  $\Sat{N}{\lforall[y][\lforall[x][(y = x + x \lor y = x + x +
        \Obj{0}')]]}$.
  પ્રાથમિક સમકક્ષતાથી
  દરેક $y \in
  \Domain{M}$ માટે
  એવો $x \in
  \Domain{M}$ છે કે
  કાં તો $x \oplus x
  = y$ અથવા
  $x \oplus x \oplus
  \mathbf{z}^*= y$.
  $x$ પ્રમાણભૂત
  હોત તો $y$
  પણ પ્રમાણભૂત
  હોત; એટલે
  $x$ અપ્રમાણભૂત છે.
  વળી $x$ અને
  $y$ જુદા ખંડોમાં
  છે, એટલે
  $x \not\approx y$.
  તે જોવા, બંને
  એક જ ખંડમાં
  હોવાનું ધારો;
  એટલે કોઈ
  પ્રમાણભૂત $n$
  માટે $x \oplus n =
  y$. જો $y = x
  \oplus x$ હોય,
  તો $x \oplus n =
  x \oplus x$;
  સરવાળાના રદ્દીકરણ
  નિયમથી $x = n$
  થાય, કારણ કે
  આ નિયમ $\Struct{N}$
  અને તેથી
  $\Struct{M}$માં
  પણ સાચો છે.
  પરંતુ ત્યારે
  $x$ પ્રમાણભૂત
  થાત. $y = x
  \oplus x \oplus
  \mathbf{z}^*$ હોય
  ત્યારે પણ
  સમાન દલીલ ચાલે છે.
\item સમાન પ્રકારની
  દલીલથી મહત્તમ
  ખંડ પણ નથી.
\item ખંડોનો ક્રમ
  સઘન છે: $[x]$
  એ $[y]$ કરતાં
  નાનો હોય અને
  $x \not\approx y$
  હોય, તો બંનેથી
  જુદો એવો વચ્ચેનો
  ખંડ~$[z]$ છે.
  ધારો કે $x \prec
  y$. અગાઉની જેમ
  $x \oplus y$ને
  બે વડે ભાગી
  શકાય, કદાચ
  શેષ સાથે;
  તેથી કોઈ $u
  \in \Domain{M}$
  માટે કાં તો
  $x \oplus y =
  u \oplus u$
  અથવા $x \oplus
  y = u \oplus u
  \oplus \mathbf{z}^*$.
  $u$ એ $x$
  અને $y$ની
  સરેરાશ છે,
  એટલે તેમની
  વચ્ચે છે. પહેલો
  કિસ્સો $x \oplus
  y = u \oplus u$
  લો; બીજો
  સમાન છે. જો
  $u \approx x$
  હોય તો કોઈ
  પ્રમાણભૂત $n$
  માટે
  \[
  x \oplus y = x \oplus n \oplus x \oplus n,
  \]
  તેથી $y = x
  \oplus n \oplus n$
  મળે અને
  $x \approx y$
  થાય, જે
  ધારણાની વિરુદ્ધ
  છે. આથી
  $u \not\approx x$.
  સમાન દલીલથી
  $u \not\approx y$.
\end{enumerate}
આથી અપ્રમાણભૂત
ખંડો સંમેય
સંખ્યાઓની જેમ
ક્રમબદ્ધ છે:
તેઓ અંતબિંદુવિહીન
!!a{enumerable}
રેખીય ક્રમ બનાવે
છે. તેથી સાચા
અંકગણિતના
કોઈ પણ બે
!!{enumerable}
અપ્રમાણભૂત
નિદર્શો
$\Struct{M}_1$
અને $\Struct{M}_2$
ના માત્ર $<$
અને $=$ ધરાવતી
ભાષા સુધીના
અવશેષો એકરૂપ
છે. ખરેખર,
એકરૂપતા~$h$
આમ વ્યાખ્યાયિત
કરી શકાય:
$\Struct{M_1}$
અને $\Struct{M}_2$
ના પ્રમાણભૂત
ભાગો પ્રમાણભૂત
નિદર્શ
$\Struct{N}$ને,
અને તેથી
એકબીજાને,
એકરૂપ છે.
અપ્રમાણભૂત
ભાગ બનાવતા
ખંડો પોતે
સંમેય સંખ્યાઓની
જેમ ક્રમબદ્ધ
છે અને
\olref[dlo]{thm:cantorQ}
મુજબ એકરૂપ
છે; ખંડોની
એકરૂપતાને
દરેક ખંડની
\emph{અંદર}ની
એકરૂપતા સુધી
વિસ્તારી શકાય:
દરેક જોડીના
ખંડોમાં મનસ્વી
ઘટકોને સામસામે
ગોઠવો અને પછી
$\Struct{M_1}$માં
$x$ના ઉત્તરગામીનું
પ્રતિબિંબ,
$\Struct{M}_2$માં
$x$ના પ્રતિબિંબના
ઉત્તરગામી તરીકે
લો. એમાંથી
$\mathfrak{M}_1$
અને $\mathfrak{M}_2$
અંકગણિતની પૂર્ણ
ભાષામાં એકરૂપ
છે એવું
\emph{ફલિત થતું નથી}.
ખરેખર, એકરૂપતા
હંમેશાં કોઈ
સંકેતમાળા સાપેક્ષ
હોય છે; $\Domain{M_1}$
અને $\Domain{M_2}$
પર સરવાળો અને
ગુણાકાર વ્યાખ્યાયિત
કરવાની પરસ્પર
બિનએકરૂપ રીતો
છે. વૉટના
પ્રસિદ્ધ પ્રમેયથી
પણ આ ફલિત
થાય છે: પૂર્ણ
સિદ્ધાંતના
ગણનીય નિદર્શોની
સંખ્યા~2
હોઈ શકે નહિ.
\sourcecorrection{OLMD-194}{મૂળમાં
બીજા નિદર્શ માટે
$\Struct{M_2}$ લખાયું
છે, જ્યારે પ્રથમ
$\Struct{M}_1$ છે;
સંગત જોડીનાં
સંજ્ઞારૂપ
$\Struct{M}_2$
અહીં વપરાયું છે.}

\begin{prob}
અંકગણિતના અપ્રમાણભૂત નિદર્શમાં
મહત્તમ ખંડ હોઈ શકે નહીં તે બતાવો.
\end{prob}

\begin{prob}
$\Lang{L}$ એ એવી પ્રથમ-ક્રમ !!{language}
હોય જેમાં $<$ એકમાત્ર !!{predicate}
છે ($\eq$ સિવાય), અને
$\Struct{N} = (\Nat,
<)$ લો. $\Struct{N}$ના બધા
સાન્ત અથવા સહસાન્ત ઉપગણો
વ્યાખ્યેય છે. આ જ
$\Struct{N}$ના \emph{એકમાત્ર}
વ્યાખ્યેય ઉપગણો છે તે
બતાવો.

(સંકેત: પહેલાં
$\Obj{prc}(x,y)$ એ
$\Lang{L}$-સૂત્ર
“$x$ એ $y$નો
તરતનો પૂર્વગામી છે”
તેનું ટૂંકું સ્વરૂપ
હોય:
\[
x<y \land \lnot \lexists[z][(x<z \land z < y)].
\]
હવે $\Struct{N}$ના
કોઈ પણ વ્યાખ્યેય
ઉપગણને અનુરૂપ
$\Lang{L}$માં
$!A(x)$ સૂત્ર છે.
એવા કોઈ પણ
$!A$ માટે
આ વાક્ય~$!D$
વિચારો:
\[
\lexists[x][\lforall[y][\lforall[z][((x<y \land x<z \land \Obj{prc}(y,z)
  \land !A(y)) \lif !A(z))]]].
\]
$\Sat{N}{!D}$ ત્યારે
અને માત્ર ત્યારે
જ્યારે $!A$ વડે
વ્યાખ્યાયિત
$\Struct{N}$નો
ઉપગણ સાન્ત
અથવા સહસાન્ત
હોય તે બતાવો.

હવે $\Struct{M}$
એ $\Struct{N}$
સાથે પ્રાથમિક
રીતે સમકક્ષ
અપ્રમાણભૂત
નિદર્શ હોય.
જો $a \in
\Domain{M}$
અપ્રમાણભૂત હોય,
તો $a$ કરતાં
મોટાં $b, c \in
\Domain{M}$ લો
અને $b$ એ~$c$નો
તરતનો પૂર્વગામી
હોય. ત્યારે
$\Domain{M}$ના
માત્ર ક્રમના
અવશેષ પર એવી
સ્વરૂપસંરક્ષક
આપલે~$h$ છે
કે $h(b)=c$;
કેમ? તેથી
$b$ એ $!A$
સંતોષે તો
$c$ પણ તેને
સંતોષે; કેમ?
આથી $!D$
એ $\Struct{M}$માં
અને તેથી
$\Struct{N}$માં
પણ સાચું છે.
પરંતુ તેનાથી
$!A$ વડે
વ્યાખ્યાયિત
$\Struct{N}$નો
ઉપગણ સાન્ત
અથવા સહસાન્ત
છે એવું ફલિત
થાય છે.
\sourcecorrection{OLMD-195}{મૂળમાં
$h$ને માત્ર વ્યાપકક્ષેત્રની
સ્વરૂપસંરક્ષક આપલે કહે છે.
અહીં તેને $<$ ધરાવતા
ક્રમ-અવશેષની આપલે
સ્પષ્ટ કહી છે; પૂર્ણ
અંકગણિતીય સંરચનાની
આપલે હોવાનો
દાવો જરૂરી નથી.}
\end{prob}

\end{document}
